\documentclass[11pt,a4paper]{article}
\usepackage{fontspec,amsmath,amssymb,amsthm,microtype}
\setmainfont{lmroman10-regular.otf}[BoldFont=lmroman10-bold.otf,ItalicFont=lmroman10-italic.otf,BoldItalicFont=lmroman10-bolditalic.otf]
\setmonofont{lmmono10-regular.otf}[Scale=MatchLowercase]
\usepackage[margin=25mm]{geometry}
\usepackage[unicode,hidelinks]{hyperref}
\setlength{\emergencystretch}{3em}
\setlength{\parskip}{0.4em}
\setlength{\parindent}{0pt}
\allowdisplaybreaks[2]
\raggedbottom
\title{\texorpdfstring{A checked Gaussian truncation bound with dependent position and velocity}{A checked Gaussian truncation bound with dependent position and velocity}}
\author{LangevinLean mathematical exposition}
\date{}
\begin{document}\maketitle
\noindent Source module: \texttt{Kinetic\allowbreak{}Langevin/\allowbreak{}Gaussian\allowbreak{}Truncation.\allowbreak{}lean}. The accompanying HTML pairs this mathematical exposition with the exact Lean statements and records the scope of the formalization.


A Hessian difference can be small through its spatial Lipschitz bound and uniformly bounded through the Hessian norm bound. This module combines those two genuine pointwise controls with the actual Gaussian velocity tail. The coupled position and velocity may be dependent; no independence premise is used.


\section{\texorpdfstring{Separate the moderate velocity region from the exact Gaussian excess}{Separate the moderate velocity region from the exact Gaussian excess}}


Suppose \(d\le\min\{B,Lx\}\,v\), with all quantities nonnegative. For a cutoff \(a\ge0\), the exact scalar estimate is \[d^2\le L^2ax^2+B^2(v^2-a)_+.\] Put \(k=\min\{B,Lx\}\). Where \(v^2\le a\), use \(k\le Lx\) and \(d^2\le k^2v^2\le L^2x^2a\). Where \(v^2>a\), write \(k^2v^2=k^2a+k^2(v^2-a)\), use \(k\le Lx\) in the first term and \(k\le B\) in the second. Lean proves the two squared bounds from the nonnegative minimum, splits on the cutoff comparison, rewrites the positive part in each region, and applies these two scalar comparisons. The subtraction of \(a\) in the tail term is retained exactly, which permits the sharp positive-part tail bound.



{\raggedright\noindent Lean declarations: \texttt{gaussian\_\allowbreak{}truncation\_\allowbreak{}pointwise}.\par}

\section{\texorpdfstring{Identify actual Hilbert L\ensuremath{{}^{2}} norms with their actual second moments}{Identify actual Hilbert L² norms with their actual second moments}}


For an actual Hilbert-valued L\ensuremath{{}^{2}} element \(X\), \(\|X\|_{L^2}^2=\int\|X(\omega)\|^2\,dP\). Lean rewrites the squared norm as the real inner product of \(X\) with itself, unfolds the actual L\ensuremath{{}^{2}} inner product as its Bochner integral, and rewrites the pointwise inner products as squared norms. This gives the precise identity needed to integrate the scalar cutoff estimate; it does not replace a possibly divergent moment by a totalized real expression.



{\raggedright\noindent Lean declarations: \texttt{lp\_\allowbreak{}norm\_\allowbreak{}sq\_\allowbreak{}eq\_\allowbreak{}integral}.\par}

\section{\texorpdfstring{Integrate the cutoff and use only the actual Gaussian velocity marginal}{Integrate the cutoff and use only the actual Gaussian velocity marginal}}


Let \(D,X,V\) be actual Euclidean L\ensuremath{{}^{2}} random variables on one probability space, and suppose \(V\) has the actual standard Gaussian law in dimension \(p\). Assume \(\|D\|\le\min\{B,L\|X\|\}\|V\|\) almost everywhere. The preceding scalar inequality integrates to \[\|D\|_{L^2}^2\le L^2a\|X\|_{L^2}^2+B^2\mathbb E(\|V\|^2-a)_+.\] Lean proves integrability of the velocity positive part directly from the actual L\ensuremath{{}^{2}} moment, proves integrability of both terms on the right, and applies almost-everywhere monotonicity of the Bochner integral.



The checked Gaussian tail theorem, transferred through the actual velocity law, gives \(\mathbb E(\|V\|^2-a)_+\le4e^{-(a-p)/4}\) for \(a\ge5p\). Hence \[\|D\|_{L^2}\le L\sqrt a\,\|X\|_{L^2}+2B e^{-(a-p)/8}.\] Lean checks the exponential-square identity, compares the two nonnegative squares, and bounds the nonnegative cross term. The only probability fact used is the marginal law of \(V\); arbitrary dependence between \(X\), \(D\), and \(V\) is retained throughout. This is the cutoff estimate required for a random Hessian-difference factor multiplying a stationary Gaussian velocity.



{\raggedright\noindent Lean declarations: \texttt{gaussian\_\allowbreak{}truncation\_\allowbreak{}Lp}.\par}

\end{document}
